\documentclass[11pt,a4paper]{article}
\usepackage[utf8]{inputenc}
\usepackage[T1]{fontenc}
\usepackage[spanish,es-minimal,es-noshorthands]{babel}
\usepackage{lmodern}
\usepackage[sfdefault,lining]{FiraSans}
\usepackage[a4paper,top=2.3cm,bottom=2.2cm,left=1.8cm,right=1.8cm,headheight=1.3cm,headsep=0.4cm,footskip=1cm]{geometry}
\usepackage{amsmath,amssymb}
\allowdisplaybreaks
\usepackage{icomma}
\usepackage{graphicx}
\usepackage[table]{xcolor}
\usepackage{tikz}
\usepackage{fontawesome5}
\usepackage{ccicons}
\usepackage{fancyhdr}
\usepackage{lastpage}
\usepackage{tasks}
\usepackage{needspace}
\usepackage[most]{tcolorbox}
\usepackage{array}
\usepackage[hidelinks]{hyperref}
\AtBeginDocument{\hypersetup{pdftitle={Soluciones de la ficha de la unidad 3: Vectores y aplicaciones. Rectas, planos y haces de planos. 2º Bach.},pdfsubject={Matemáticas II 2º Bach.},pdfkeywords={vectores, producto escalar, producto vectorial, producto mixto, áreas, volúmenes, ecuaciones de la recta, ecuaciones del plano, posiciones relativas, haces de planos, PAU}}}

\definecolor{navy}{HTML}{1B2748}
\definecolor{accent}{HTML}{E8603A}
\definecolor{teal}{HTML}{1F9E90}
\definecolor{soft}{HTML}{EEF1F8}
\definecolor{muted}{HTML}{5D6685}

\setlength{\parindent}{0pt}
\setlength{\parskip}{3pt}
\settasks{label=\alph*), label-format=\bfseries\color{navy}, label-width=1.3em, item-indent=1.8em, column-sep=0.6em, after-item-skip=5pt, before-skip=3pt, after-skip=4pt}

% ---------- capçalera i peu ----------
\newcommand{\llicencia}{\href{https://creativecommons.org/licenses/by-nc-sa/4.0/deed.es}{\ccbyncsa}\ \ Licencia Creative Commons Atribución-NoComercial-CompartirIgual 4.0 Internacional (CC BY-NC-SA 4.0)}
\pagestyle{fancy}
\fancyhf{}
\fancyhead[L]{\raisebox{-0.35\height}{\includegraphics[height=0.95cm]{logo-gva.png}}}
\fancyhead[C]{\small\color{muted}Soluciones de la ficha · Unidad 3: Vectores, rectas y planos · 2\textsuperscript{o} Bach.}
\fancyhead[R]{\raisebox{-0.35\height}{\includegraphics[height=0.95cm]{logo-vermellar.png}}}
\fancyfoot[L]{\scriptsize\color{muted}\llicencia}
\fancyfoot[R]{\small\color{muted}\thepage}
\renewcommand{\headrulewidth}{0.4pt}
\renewcommand{\headrule}{\color{muted!40}\hrule height \headrulewidth}
\fancypagestyle{primera}{\fancyhf{}\renewcommand{\headrulewidth}{0pt}
  \fancyfoot[L]{\scriptsize\color{muted}\llicencia}
  \fancyfoot[R]{\small\color{muted}\thepage}}

% ---------- nivells i exercicis ----------
\newcommand{\estrelles}[1]{{\color{accent}\ifcase#1\or\faStar\faStar[regular]\faStar[regular]\or\faStar\faStar\faStar[regular]\or\faStar\faStar\faStar\fi}}
\newcounter{exer}
\newcommand{\ex}[1]{\par\needspace{4\baselineskip}\medskip\refstepcounter{exer}\noindent{\color{navy}\bfseries\theexer.}\ \ }

\newcommand{\seccio}[2]{\par\needspace{10\baselineskip}\bigskip
  \begin{tcolorbox}[colback=navy,colframe=navy,coltext=white,arc=2mm,boxrule=0pt,left=3mm,right=3mm,top=1.2mm,bottom=1.2mm]
  {\large\bfseries #1\quad #2}
  \end{tcolorbox}\vspace{-2pt}}


\setlength{\emergencystretch}{3em}  % les fórmules llargues (equacions de rectes, ternes) passen a la línia següent

% quadre de recordatori (només a la fitxa d'exercicis)
\newcommand{\recorda}[1]{\par\smallskip\begin{tcolorbox}[colback=accent!7,colframe=accent!55,arc=2mm,boxrule=0.5pt,left=2.5mm,right=2.5mm,top=1mm,bottom=1mm]\small{\color{accent}\faLightbulb[regular]}\ \textbf{Recuerda:}\ #1\end{tcolorbox}}
% etiqueta dels exercicis de la PAU d'altres comunitats
\newcommand{\pau}[1]{{\color{teal}\small\bfseries\faUniversity\ PAU · #1.}\ }
\newcommand{\res}[1]{\boxed{#1}}\newcommand{\rest}[1]{\fbox{#1}}

% =====================================================================
% Capçalera FentMates de les fitxes: un color per a cada curs.
% Els colors són els de la pàgina inicial (LEVELS a assets/js/materials.js);
% el logo ve de _eines/logo/pdf/ (python3 _eines/logo/genera.py).
%
% Ús, al final del preàmbul de la unitat:
%   % =====================================================================
% Capçalera FentMates de les fitxes: un color per a cada curs.
% Els colors són els de la pàgina inicial (LEVELS a assets/js/materials.js);
% el logo ve de _eines/logo/pdf/ (python3 _eines/logo/genera.py).
%
% Ús, al final del preàmbul de la unitat:
%   % =====================================================================
% Capçalera FentMates de les fitxes: un color per a cada curs.
% Els colors són els de la pàgina inicial (LEVELS a assets/js/materials.js);
% el logo ve de _eines/logo/pdf/ (python3 _eines/logo/genera.py).
%
% Ús, al final del preàmbul de la unitat:
%   \input{../comu/fentmates.tex}
%   \fmcurs{1eso}
% i, al començament del document:
%   \capcalera{Matemàtiques}{1r ESO}{Fitxa d'exercicis · Unitat 2}{Divisibilitat}
% Opcional abans de \input: \def\fmneed{14} (espai mínim abans de cada secció).
% =====================================================================
\makeatletter
\def\fm@logodir{../../logo/pdf}
\providecommand{\fmneed}{10}
\def\fm@cursnom{}

% clau del curs, color i color del text damunt del color
\newcommand{\fm@curs}[3]{\@namedef{fm@color@#1}{#2}\@namedef{fm@text@#1}{#3}}
\fm@curs{1eso}{E6B84A}{navy}
\fm@curs{2eso}{3B82F6}{white}
\fm@curs{3eso}{FF7A50}{white}
\fm@curs{4eso-a}{A78BFA}{white}
\fm@curs{4eso-b}{6D28D9}{white}
\fm@curs{1batx-mi}{9B1C1C}{white}
\fm@curs{1batx-ccss}{EF6B6B}{white}
\fm@curs{2batx-mii}{1E6B47}{white}
\fm@curs{2batx-ccss}{2BB3A3}{white}

\newcommand{\fmcurs}[1]{%
  \@ifundefined{fm@color@#1}{\PackageError{fentmates}{Curs desconegut: #1}{}}{}%
  \edef\fm@tmp{\noexpand\definecolor{curs}{HTML}{\@nameuse{fm@color@#1}}%
               \noexpand\colorlet{curstext}{\@nameuse{fm@text@#1}}}\fm@tmp
  \colorlet{cursfosc}{curs!72!black}%
  \def\fm@logo{\fm@logodir/vertical-#1.pdf}%
  \def\fm@marca{\fm@logodir/marca-#1.pdf}}

% ---------- capçalera de la primera pàgina ----------
% #1 matèria (petit), #2 curs, #3 tipus de document i unitat, #4 títol
\newcommand{\capcalera}[4]{%
  \gdef\fm@cursnom{#2}%
  \par\noindent
  \begin{tcolorbox}[colback=curs,colframe=curs,boxrule=0pt,arc=2.2mm,width=3.5cm,height=2.15cm,
      valign=center,left=3mm,right=2mm,top=1mm,bottom=1mm,nobeforeafter,box align=center]
    \color{curstext}{\footnotesize #1}\par\vspace{1pt}{\fontsize{19}{21}\selectfont\bfseries #2\par}%
  \end{tcolorbox}\hspace{5mm}%
  \begin{minipage}[c]{\dimexpr\linewidth-3.5cm-5mm-2.9cm\relax}
    {\color{cursfosc}\small\bfseries\MakeUppercase{#3}}\par\vspace{2pt}
    {\color{navy}\fontsize{20}{23}\selectfont\bfseries #4\par}
  \end{minipage}\hfill
  \begin{minipage}[c]{2.6cm}\raggedleft\includegraphics[width=2.6cm]{\fm@logo}\end{minipage}\par
  \vspace{2pt}}

% ---------- barres de secció del color del curs ----------
\renewcommand{\seccio}[2]{\par\needspace{\fmneed\baselineskip}\bigskip
  \begin{tcolorbox}[colback=curs,colframe=curs,coltext=curstext,arc=2mm,boxrule=0pt,left=3mm,right=3mm,top=1.2mm,bottom=1.2mm]
  {\large\bfseries #1\quad #2}
  \end{tcolorbox}\vspace{-2pt}}

% ---------- estrelles de dificultat del color del curs (a Batxillerat no se'n posen) ----------
\@ifundefined{estrelles}{}{%
  \renewcommand{\estrelles}[1]{{\color{curs}\ifcase#1\or\faStar\faStar[regular]\faStar[regular]\or\faStar\faStar\faStar[regular]\or\faStar\faStar\faStar\fi}}}

% ---------- capçalera i peu de les pàgines ----------
% primera pàgina: logos institucionals al peu
\fancypagestyle{primera}{\fancyhf{}\renewcommand{\headrulewidth}{0pt}
  \fancyfoot[L]{\raisebox{-0.32\height}{\includegraphics[height=0.78cm]{logo-gva.png}}\hspace{5mm}%
    \raisebox{-0.32\height}{\includegraphics[height=0.78cm]{logo-vermellar.png}}}
  \fancyfoot[R]{\scriptsize\color{muted}\href{https://creativecommons.org/licenses/by-nc-sa/4.0/deed.ca}{\ccbyncsa}\ \ CC BY-NC-SA 4.0 · fentmates.org\qquad\small\thepage}}
% la resta: el símbol a l'esquerra i el curs a la dreta, amb el títol al centre com abans
\fancyhead[L]{\raisebox{-0.3\height}{\includegraphics[height=0.62cm]{\fm@marca}}}
\fancyhead[R]{\small\bfseries\color{cursfosc}\fm@cursnom}
\renewcommand{\headrule}{\color{curs!45}\hrule height \headrulewidth}
\makeatother

%   \fmcurs{1eso}
% i, al començament del document:
%   \capcalera{Matemàtiques}{1r ESO}{Fitxa d'exercicis · Unitat 2}{Divisibilitat}
% Opcional abans de \input: \def\fmneed{14} (espai mínim abans de cada secció).
% =====================================================================
\makeatletter
\def\fm@logodir{../../logo/pdf}
\providecommand{\fmneed}{10}
\def\fm@cursnom{}

% clau del curs, color i color del text damunt del color
\newcommand{\fm@curs}[3]{\@namedef{fm@color@#1}{#2}\@namedef{fm@text@#1}{#3}}
\fm@curs{1eso}{E6B84A}{navy}
\fm@curs{2eso}{3B82F6}{white}
\fm@curs{3eso}{FF7A50}{white}
\fm@curs{4eso-a}{A78BFA}{white}
\fm@curs{4eso-b}{6D28D9}{white}
\fm@curs{1batx-mi}{9B1C1C}{white}
\fm@curs{1batx-ccss}{EF6B6B}{white}
\fm@curs{2batx-mii}{1E6B47}{white}
\fm@curs{2batx-ccss}{2BB3A3}{white}

\newcommand{\fmcurs}[1]{%
  \@ifundefined{fm@color@#1}{\PackageError{fentmates}{Curs desconegut: #1}{}}{}%
  \edef\fm@tmp{\noexpand\definecolor{curs}{HTML}{\@nameuse{fm@color@#1}}%
               \noexpand\colorlet{curstext}{\@nameuse{fm@text@#1}}}\fm@tmp
  \colorlet{cursfosc}{curs!72!black}%
  \def\fm@logo{\fm@logodir/vertical-#1.pdf}%
  \def\fm@marca{\fm@logodir/marca-#1.pdf}}

% ---------- capçalera de la primera pàgina ----------
% #1 matèria (petit), #2 curs, #3 tipus de document i unitat, #4 títol
\newcommand{\capcalera}[4]{%
  \gdef\fm@cursnom{#2}%
  \par\noindent
  \begin{tcolorbox}[colback=curs,colframe=curs,boxrule=0pt,arc=2.2mm,width=3.5cm,height=2.15cm,
      valign=center,left=3mm,right=2mm,top=1mm,bottom=1mm,nobeforeafter,box align=center]
    \color{curstext}{\footnotesize #1}\par\vspace{1pt}{\fontsize{19}{21}\selectfont\bfseries #2\par}%
  \end{tcolorbox}\hspace{5mm}%
  \begin{minipage}[c]{\dimexpr\linewidth-3.5cm-5mm-2.9cm\relax}
    {\color{cursfosc}\small\bfseries\MakeUppercase{#3}}\par\vspace{2pt}
    {\color{navy}\fontsize{20}{23}\selectfont\bfseries #4\par}
  \end{minipage}\hfill
  \begin{minipage}[c]{2.6cm}\raggedleft\includegraphics[width=2.6cm]{\fm@logo}\end{minipage}\par
  \vspace{2pt}}

% ---------- barres de secció del color del curs ----------
\renewcommand{\seccio}[2]{\par\needspace{\fmneed\baselineskip}\bigskip
  \begin{tcolorbox}[colback=curs,colframe=curs,coltext=curstext,arc=2mm,boxrule=0pt,left=3mm,right=3mm,top=1.2mm,bottom=1.2mm]
  {\large\bfseries #1\quad #2}
  \end{tcolorbox}\vspace{-2pt}}

% ---------- estrelles de dificultat del color del curs (a Batxillerat no se'n posen) ----------
\@ifundefined{estrelles}{}{%
  \renewcommand{\estrelles}[1]{{\color{curs}\ifcase#1\or\faStar\faStar[regular]\faStar[regular]\or\faStar\faStar\faStar[regular]\or\faStar\faStar\faStar\fi}}}

% ---------- capçalera i peu de les pàgines ----------
% primera pàgina: logos institucionals al peu
\fancypagestyle{primera}{\fancyhf{}\renewcommand{\headrulewidth}{0pt}
  \fancyfoot[L]{\raisebox{-0.32\height}{\includegraphics[height=0.78cm]{logo-gva.png}}\hspace{5mm}%
    \raisebox{-0.32\height}{\includegraphics[height=0.78cm]{logo-vermellar.png}}}
  \fancyfoot[R]{\scriptsize\color{muted}\href{https://creativecommons.org/licenses/by-nc-sa/4.0/deed.ca}{\ccbyncsa}\ \ CC BY-NC-SA 4.0 · fentmates.org\qquad\small\thepage}}
% la resta: el símbol a l'esquerra i el curs a la dreta, amb el títol al centre com abans
\fancyhead[L]{\raisebox{-0.3\height}{\includegraphics[height=0.62cm]{\fm@marca}}}
\fancyhead[R]{\small\bfseries\color{cursfosc}\fm@cursnom}
\renewcommand{\headrule}{\color{curs!45}\hrule height \headrulewidth}
\makeatother

%   \fmcurs{1eso}
% i, al començament del document:
%   \capcalera{Matemàtiques}{1r ESO}{Fitxa d'exercicis · Unitat 2}{Divisibilitat}
% Opcional abans de \input: \def\fmneed{14} (espai mínim abans de cada secció).
% =====================================================================
\makeatletter
\def\fm@logodir{../../logo/pdf}
\providecommand{\fmneed}{10}
\def\fm@cursnom{}

% clau del curs, color i color del text damunt del color
\newcommand{\fm@curs}[3]{\@namedef{fm@color@#1}{#2}\@namedef{fm@text@#1}{#3}}
\fm@curs{1eso}{E6B84A}{navy}
\fm@curs{2eso}{3B82F6}{white}
\fm@curs{3eso}{FF7A50}{white}
\fm@curs{4eso-a}{A78BFA}{white}
\fm@curs{4eso-b}{6D28D9}{white}
\fm@curs{1batx-mi}{9B1C1C}{white}
\fm@curs{1batx-ccss}{EF6B6B}{white}
\fm@curs{2batx-mii}{1E6B47}{white}
\fm@curs{2batx-ccss}{2BB3A3}{white}

\newcommand{\fmcurs}[1]{%
  \@ifundefined{fm@color@#1}{\PackageError{fentmates}{Curs desconegut: #1}{}}{}%
  \edef\fm@tmp{\noexpand\definecolor{curs}{HTML}{\@nameuse{fm@color@#1}}%
               \noexpand\colorlet{curstext}{\@nameuse{fm@text@#1}}}\fm@tmp
  \colorlet{cursfosc}{curs!72!black}%
  \def\fm@logo{\fm@logodir/vertical-#1.pdf}%
  \def\fm@marca{\fm@logodir/marca-#1.pdf}}

% ---------- capçalera de la primera pàgina ----------
% #1 matèria (petit), #2 curs, #3 tipus de document i unitat, #4 títol
\newcommand{\capcalera}[4]{%
  \gdef\fm@cursnom{#2}%
  \par\noindent
  \begin{tcolorbox}[colback=curs,colframe=curs,boxrule=0pt,arc=2.2mm,width=3.5cm,height=2.15cm,
      valign=center,left=3mm,right=2mm,top=1mm,bottom=1mm,nobeforeafter,box align=center]
    \color{curstext}{\footnotesize #1}\par\vspace{1pt}{\fontsize{19}{21}\selectfont\bfseries #2\par}%
  \end{tcolorbox}\hspace{5mm}%
  \begin{minipage}[c]{\dimexpr\linewidth-3.5cm-5mm-2.9cm\relax}
    {\color{cursfosc}\small\bfseries\MakeUppercase{#3}}\par\vspace{2pt}
    {\color{navy}\fontsize{20}{23}\selectfont\bfseries #4\par}
  \end{minipage}\hfill
  \begin{minipage}[c]{2.6cm}\raggedleft\includegraphics[width=2.6cm]{\fm@logo}\end{minipage}\par
  \vspace{2pt}}

% ---------- barres de secció del color del curs ----------
\renewcommand{\seccio}[2]{\par\needspace{\fmneed\baselineskip}\bigskip
  \begin{tcolorbox}[colback=curs,colframe=curs,coltext=curstext,arc=2mm,boxrule=0pt,left=3mm,right=3mm,top=1.2mm,bottom=1.2mm]
  {\large\bfseries #1\quad #2}
  \end{tcolorbox}\vspace{-2pt}}

% ---------- estrelles de dificultat del color del curs (a Batxillerat no se'n posen) ----------
\@ifundefined{estrelles}{}{%
  \renewcommand{\estrelles}[1]{{\color{curs}\ifcase#1\or\faStar\faStar[regular]\faStar[regular]\or\faStar\faStar\faStar[regular]\or\faStar\faStar\faStar\fi}}}

% ---------- capçalera i peu de les pàgines ----------
% primera pàgina: logos institucionals al peu
\fancypagestyle{primera}{\fancyhf{}\renewcommand{\headrulewidth}{0pt}
  \fancyfoot[L]{\raisebox{-0.32\height}{\includegraphics[height=0.78cm]{logo-gva.png}}\hspace{5mm}%
    \raisebox{-0.32\height}{\includegraphics[height=0.78cm]{logo-vermellar.png}}}
  \fancyfoot[R]{\scriptsize\color{muted}\href{https://creativecommons.org/licenses/by-nc-sa/4.0/deed.ca}{\ccbyncsa}\ \ CC BY-NC-SA 4.0 · fentmates.org\qquad\small\thepage}}
% la resta: el símbol a l'esquerra i el curs a la dreta, amb el títol al centre com abans
\fancyhead[L]{\raisebox{-0.3\height}{\includegraphics[height=0.62cm]{\fm@marca}}}
\fancyhead[R]{\small\bfseries\color{cursfosc}\fm@cursnom}
\renewcommand{\headrule}{\color{curs!45}\hrule height \headrulewidth}
\makeatother

\fmcurs{2batx-mii}

\begin{document}
\thispagestyle{primera}

\capcalera{Matemáticas II}{2\textsuperscript{o} Bach.}{Soluciones de la ficha · Unidad 3}{\raggedright Vectores y aplicaciones. Rectas, planos y haces de planos}
\vspace{8pt}

\begin{tcolorbox}[colback=soft,colframe=soft,arc=2mm,left=3mm,right=3mm,top=1.5mm,bottom=1.5mm]
\small Intenta resolver cada ejercicio antes de mirar su solución. Una recta o un plano se pueden escribir de muchas maneras: si tu punto o tu vector director son distintos, comprueba que el punto es de la recta (o del plano) y que los vectores son proporcionales. Puedes comprobar los resultados con las herramientas interactivas de los apuntes.
\end{tcolorbox}

% =====================================================
\seccio{1}{Vectores en el espacio y aplicaciones}

\ex{1} \begin{tasks}(1)
  \task $\overrightarrow{AB}=(2,\,1,\,-2)$, $|\overrightarrow{AB}|=\sqrt{9}=\res{3}$; \ $\overrightarrow{AC}=(-2,\,4,\,4)$, $|\overrightarrow{AC}|=\sqrt{36}=\res{6}$
  \task \mbox{$\dfrac{\overrightarrow{AC}}{6}=\res{\left(-\tfrac13,\,\tfrac23,\,\tfrac23\right)}$}
  \task $-3\,\overrightarrow{AB}=\res{(-6,\,-3,\,6)}$
  \task $\res{(4,\,7,\,-2)}$
\end{tasks}

\ex{1} $B=2M-A=\res{(-1,\,-4,\,9)}$. \qquad $P'=2Q-P=\res{(-2,\,5,\,-5)}$

\ex{2} $\overrightarrow{AB}=(6,\,-6,\,-6)$, \ $\tfrac13\overrightarrow{AB}=(2,\,-2,\,-2)$: \ $P_1=A+\tfrac13\overrightarrow{AB}=\res{(0,\,-1,\,3)}$, \ $P_2=A+\tfrac23\overrightarrow{AB}=\res{(2,\,-3,\,1)}$

\ex{2} $\overrightarrow{AB}=(1,\,1,\,-2)$ y $\overrightarrow{AC}=(a-1,\,3,\,b-2)$ deben ser proporcionales: \mbox{$\dfrac{a-1}{1}=\dfrac31=\dfrac{b-2}{-2}$} \ $\Rightarrow$ \ $\res{a=4,\ b=-4}$

\ex{2} $\overrightarrow{AD}=\overrightarrow{BC}$: \ $D=A+C-B=\res{(0,\,1,\,1)}$. Las diagonales se cortan en el punto medio: \ centro \mbox{$=\dfrac{A+C}{2}=\res{\left(\tfrac32,\,\tfrac12,\,\tfrac32\right)}$}

\ex{2} \mbox{$G=\dfrac{A+B+C}{3}\Rightarrow C=3G-A-B=(3-2+1,\ 6-1-4,\ 0+3-2)=\res{(2,\,1,\,1)}$}

\ex{2} $P$ está a dos tercios del segmento desde $A$: \ $P=A+\tfrac23\overrightarrow{AB}=(1,\,-1,\,2)+\tfrac23(6,\,6,\,-6)=\res{(5,\,3,\,-2)}$

\ex{2} \pau{Castilla-La Mancha, extraordinaria 2026}$\overrightarrow{AB}=(2,\,4,\,2)\parallel(1,\,2,\,1)$ y $\overrightarrow{AC}=(m-1,\,m-2,\,m-1)$: hace falta $m-2=2(m-1)$, \ $\res{m=0}$ (y entonces $\overrightarrow{AC}=(-1,\,-2,\,-1)$ \checkmark).

\ex{2} \pau{Cataluña, extraordinaria 2026}La tapa es el paralelogramo de lados $AB$ y $AC$: \ $D=B+C-A=\res{(4,\,10,\,5)}$. \ $|\overrightarrow{AB}|=|(3,\,0,\,4)|=\res{5}$, \ $|\overrightarrow{AC}|=|(0,\,9,\,0)|=\res{9}$. (Además, $\overrightarrow{AB}\cdot\overrightarrow{AC}=0$: la tapa es un rectángulo.)


% =====================================================
\seccio{2}{Dependencia lineal y bases}

\ex{1} $\begin{vmatrix}1 & 2 & -1\\2 & 1 & 3\\0 & 3 & -5\end{vmatrix}=1(-5-9)-2(-10-0)-1(6-0)=0$: \ \rest{linealmente dependientes}. \ $\vec w=a\vec u+b\vec v$: \ $a+2b=0$, $2a+b=3$ $\Rightarrow$ $a=2$, $b=-1$, y la tercera se cumple ($-2-3=-5$): \ $\res{\vec w=2\vec u-\vec v}$

\ex{2} \begin{tasks}(1)
  \task $\begin{vmatrix}1&1&1\\1&a&1\\1&1&a^2\end{vmatrix}=a^3-a^2-a+1=(a-1)^2(a+1)$: \ base si \ $\res{a\neq1\ \text{y}\ a\neq-1}$
  \task Si $a=1$, los tres vectores son $(1,\,1,\,1)$: \ rango $\res{1}$. Si $a=-1$, $\vec w=\vec u$ pero $\vec v=(1,\,-1,\,1)$ no es proporcional a $\vec u$: \ rango $\res{2}$.
\end{tasks}

\ex{2} El determinante vale $2\neq0$: es una base. \ $a(1,\,1,\,0)+b(0,\,1,\,1)+c(1,\,0,\,1)=(3,\,1,\,2)$: \ $a+c=3$, $a+b=1$, $b+c=2$. Sumando, $a+b+c=3$, de donde $c=2$, $a=1$, $b=0$: \ $\res{\vec x=(1,\,0,\,2)_{\mathcal B}}$

\ex{2} $a(1,\,1,\,0)+b(0,\,1,\,1)=(a,\,a+b,\,b)$. Para $(1,\,2,\,3)$: $a=1$, $b=3$, pero $a+b=4\neq2$: \ \rest{no}. Para $(2,\,5,\,3)$: $a=2$, $b=3$ y $a+b=5$ \ \checkmark: \ $\res{(2,\,5,\,3)=2(1,\,1,\,0)+3(0,\,1,\,1)}$

\ex{3} \begin{tasks}(1)
  \task Coordenadas $(1,\,1,\,0)$, $(0,\,1,\,1)$, $(1,\,0,\,1)$: determinante $2\neq0$: \ \rest{independientes}.
  \task La suma de los tres vectores es $\vec 0$ (el determinante de las coordenadas es 0): \ \rest{dependientes}.
\end{tasks}


% =====================================================
\seccio{3}{Producto escalar}

\ex{1} \begin{tasks}(1)
  \task $\vec u\cdot\vec v=3-8=\res{-5}$; $|\vec u|=\res{3}$, $|\vec v|=\res{5}$
  \task \mbox{$\cos\alpha=\dfrac{-5}{15}=-\dfrac13\Rightarrow\res{\alpha\approx109^\circ28'16''}$}
  \task \mbox{$\dfrac{\vec u\cdot\vec v}{|\vec v|}=\res{-1}$}
  \task \mbox{$\dfrac{\vec u\cdot\vec v}{|\vec v|^2}\vec v=-\dfrac15(3,\,0,\,-4)=\res{\left(-\tfrac35,\,0,\,\tfrac45\right)}$}
\end{tasks}

\ex{2} $(x,\,y,\,z)$: \ $x+2y-z=0$, \ $2x+z=0$ $\Rightarrow$ $z=-2x$, $y=-\tfrac32x$: los vectores son proporcionales a $(2,\,-3,\,-4)$, que tiene módulo $\sqrt{29}$: \ $\res{\pm(2,\,-3,\,-4)}$

\ex{2} \mbox{$\cos60^\circ=\dfrac{4+m}{\sqrt{5+m^2}\,\sqrt6}=\dfrac12\Rightarrow2(4+m)=\sqrt6\sqrt{5+m^2}$}. Elevando al cuadrado: $m^2-16m-17=0$, \ $m=17$ o $m=-1$. Ambas hacen $4+m>0$: \ $\res{m=17\ \text{o}\ m=-1}$

\ex{2} $\cos\alpha=\tfrac23$, $\cos\beta=-\tfrac13$, $\cos\gamma=\tfrac23$: \ $\res{\alpha=\gamma\approx48^\circ11',\ \beta\approx109^\circ28'}$. Como $\cos^2\alpha+\cos^2\beta+\cos^2\gamma=1$: $\tfrac12+\tfrac14+\cos^2\gamma=1\Rightarrow\cos\gamma=\pm\tfrac12$: \ sí, y $\res{\gamma=60^\circ\ \text{o}\ 120^\circ}$

\ex{3} \begin{tasks}(1)
  \task $\cos\alpha=\tfrac{6}{12}=\tfrac12$: \ $\res{60^\circ}$
  \task $|\vec u+\vec v|^2=9+16+12=37$: $\res{\sqrt{37}}$; \ $|\vec u-\vec v|^2=9+16-12$: $\res{\sqrt{13}}$
  \task $2\cdot9+6\cdot6-6-3\cdot16=\res{0}$: \ los vectores $2\vec u-\vec v$ y $\vec u+3\vec v$ son perpendiculares.
\end{tasks}

\ex{2} $\overrightarrow{AB}=(2,\,1,\,-2)$, $\overrightarrow{AC}=(1,\,2,\,2)$: $\overrightarrow{AB}\cdot\overrightarrow{AC}=0$ y $|\overrightarrow{AB}|=|\overrightarrow{AC}|=3$: rectángulo en $A$ e isósceles. \ Ángulos: \ $\res{\hat A=90^\circ,\ \hat B=\hat C=45^\circ}$

\ex{1} $2m+m-3=0\Rightarrow\res{m=1}$


% =====================================================
\seccio{4}{Producto vectorial y producto mixto. Áreas y volúmenes}

\ex{1} $\vec u\times\vec v=\begin{vmatrix}\vec i&\vec j&\vec k\\2&-1&3\\1&0&-2\end{vmatrix}=\res{(2,\,7,\,1)}$, \ $\vec v\times\vec u=\res{(-2,\,-7,\,-1)}$. \ $(2,\,7,\,1)\cdot(2,\,-1,\,3)=4-7+3=0$, \ $(2,\,7,\,1)\cdot(1,\,0,\,-2)=2-2=0$.

\ex{2} $\overrightarrow{AB}\times\overrightarrow{AC}=(1,\,1,\,2)\times(-1,\,2,\,0)=(-4,\,-2,\,3)$, de módulo $\sqrt{29}$: \ $\res{\pm\tfrac{1}{\sqrt{29}}(-4,\,-2,\,3)}$; \ área $=\res{\tfrac{\sqrt{29}}{2}\ u^2}$

\ex{2} $C=B+D-A=\res{(2,\,4,\,3)}$. \ $\overrightarrow{AB}\times\overrightarrow{AD}=(2,\,1,\,1)\times(-1,\,2,\,2)=(0,\,-5,\,5)$: \ área $=\sqrt{50}=\res{5\sqrt2\ u^2}$

\ex{2} $\overrightarrow{AB}\times\overrightarrow{AC}=(-1,\,1,\,0)\times(-1,\,0,\,k)=(k,\,k,\,1)$; \ $\tfrac12\sqrt{2k^2+1}=\tfrac32\Rightarrow2k^2+1=9$: \ $\res{k=\pm2}$

\ex{2} $\big[\overrightarrow{AB},\overrightarrow{AC},\overrightarrow{AD}\big]=\begin{vmatrix}2 & 0 & -1\\0 & 3 & 1\\0 & a-1 & 2\end{vmatrix}=2(6-a+1)=2(7-a)$. \ \mbox{$\dfrac{|2(7-a)|}{6}=2\Rightarrow|7-a|=6$}: \ $\res{a=1\ \text{o}\ a=13}$

\ex{2} $\overrightarrow{AB}=(1,\,1,\,-1)$, $\overrightarrow{AC}=(0,\,-2,\,2)$, $\overrightarrow{AD}=(a-1,\,-1,\,a-1)$: \ $\big[\overrightarrow{AB},\overrightarrow{AC},\overrightarrow{AD}\big]=-2a+4=0\Rightarrow\res{a=2}$

\ex{3} \begin{tasks}(1)
  \task $\res{-5}$ (hemos permutado dos filas)
  \task $2\cdot3\,[\vec u,\vec v-\vec w,\vec w]=6\,[\vec u,\vec v,\vec w]=\res{30}$
  \task $[\vec u,\vec v,\vec w]+[\vec v,\vec v,\vec w]=\res{5}$
\end{tasks}

\ex{2} \pau{Andalucía, ordinaria 2026}\begin{tasks}(1)
  \task $\big[\overrightarrow{AB},\overrightarrow{AC},\overrightarrow{AD}\big]=\begin{vmatrix}1 & m-2 & 1\\2 & 2 & 2\\0 & -3 & m\end{vmatrix}=-2m^2+6m=-2m(m-3)$: \ $\res{m=0\ \text{o}\ m=3}$
  \task $\overrightarrow{AB}\times\overrightarrow{AC}=(1,\,-1,\,1)\times(2,\,2,\,2)=(-4,\,0,\,4)$: \ área $=\tfrac12\sqrt{32}=\res{2\sqrt2\ u^2}$. \ El producto mixto vale $-2+6=4$: \ $V=\tfrac46=\res{\tfrac23\ u^3}$
\end{tasks}

\ex{2} \pau{Aragón, ordinaria 2024}\begin{tasks}(1)
  \task $\big[\overrightarrow{P_1P_2},\overrightarrow{P_1P_3},\overrightarrow{P_1P_4}\big]=\begin{vmatrix}2 & a-1 & -3\\1 & 0 & -2\\3 & 2 & -4\end{vmatrix}=4-2a$; \ \mbox{$\dfrac{|4-2a|}{6}=\dfrac13$}: \ $\res{a=1\ \text{o}\ a=3}$
  \task $4-2a=0$: \ $\res{a=2}$
\end{tasks}

\ex{2} \pau{Asturias, ordinaria 2025}\begin{tasks}(1)
  \task $\big[\overrightarrow{AB},\overrightarrow{AC},\overrightarrow{AD}\big]=\begin{vmatrix}-3 & -1 & 1\\2 & -2 & 2\\-1 & -3 & -1\end{vmatrix}=-32\neq0$: \ \rest{no son coplanarios}.
  \task $\overrightarrow{AB}\times\overrightarrow{AC}=(0,\,8,\,8)$: \ área $=\tfrac12\cdot8\sqrt2=\res{4\sqrt2\ u^2}$
  \task $V=\tfrac{32}{6}=\res{\tfrac{16}{3}\ u^3}$
\end{tasks}

\ex{2} \pau{Castilla-La Mancha, ordinaria 2026}\begin{tasks}(1)
  \task $[\vec u,\vec v,\vec w]=\begin{vmatrix}1 & -1 & a\\2 & a & 1\\-1 & 1 & -1\end{vmatrix}=a^2+a-2=(a+2)(a-1)$: \ $\res{a=-2\ \text{o}\ a=1}$
  \task $|a^2+a-2|=10$: $a^2+a-12=0$ da $\res{a=3\ \text{o}\ a=-4}$, y $a^2+a+8=0$ no tiene soluciones reales.
\end{tasks}

\ex{2} \pau{Madrid, extraordinaria 2025}\begin{tasks}(1)
  \task $\big[\overrightarrow{AB},\overrightarrow{AC},\overrightarrow{AD}\big]=\begin{vmatrix}1 & -2 & -2\\-3 & -1 & 1\\1 & -4 & -3\end{vmatrix}=-3\neq0$: no son coplanarios. \ $V=\res{\tfrac12\ u^3}$
  \task $\overrightarrow{AB}\times\overrightarrow{AC}=(-4,\,5,\,-7)$: \ $S=\res{\tfrac{3\sqrt{10}}{2}\ u^2}$. \ \mbox{$V=\tfrac13Sh\Rightarrow h=\dfrac{3V}{S}=\dfrac{3/2}{3\sqrt{10}/2}=\res{\tfrac{\sqrt{10}}{10}\ u}$}
\end{tasks}

\ex{2} \pau{Navarra, ordinaria 2024}\begin{tasks}(1)
  \task $|\overrightarrow{AB}|=|\overrightarrow{AC}|=\sqrt{21}$ y $|\overrightarrow{BC}|=\sqrt{12}$: los lados iguales son $AB$ y $AC$, y $A$ es el vértice común. \ $D=B+C-A=\res{(-2,\,-2,\,3)}$
  \task $\overrightarrow{AB}\times\overrightarrow{AC}=(-2,\,1,\,4)\times(-4,\,-1,\,2)=(6,\,-12,\,6)$: \ área $=\sqrt{216}=\res{6\sqrt6\ u^2}$
\end{tasks}

\ex{2} \pau{Aragón, extraordinaria 2024}\begin{tasks}(1)
  \task $\vec v_1+\vec v_2+\vec v_3=\overrightarrow{AA}=\vec 0$: \ \rest{son linealmente dependientes} (sin hacer ningún cálculo).
  \task $\overrightarrow{AB}\times\overrightarrow{AC}=(0,\,-2,\,-4)\times(1,\,0,\,-1)=(2,\,-4,\,2)$: \ área $=\tfrac12\sqrt{24}=\res{\sqrt6\ u^2}$
\end{tasks}


% =====================================================
\seccio{5}{Ecuaciones de rectas y planos}

\ex{1} $\vec d=\overrightarrow{AB}=(2,\,1,\,-2)$. \ $(x,\,y,\,z)=(2,\,-1,\,3)+\lambda(2,\,1,\,-2)$; \quad $\left\{\begin{aligned}x&=2+2\lambda\\y&=-1+\lambda\\z&=3-2\lambda\end{aligned}\right.$; \quad \mbox{$\dfrac{x-2}{2}=\dfrac{y+1}{1}=\dfrac{z-3}{-2}$}; \quad $\left\{\begin{aligned}x-2y-4&=0\\2y+z-1&=0\end{aligned}\right.$. \ $C$: $\lambda=-1$ da $(0,\,-2,\,5)$: \rest{$C\in r$}. \ $D$: de $x=6$ sale $\lambda=2$, pero entonces $z=-1\neq0$: \rest{$D\notin r$}.

\ex{1} $\vec d=(1,\,-1,\,1)\times(2,\,1,\,-1)=(0,\,3,\,3)\parallel(0,\,1,\,1)$. Sumando las ecuaciones: $3x=3$, $x=1$; con $z=0$, $y=-1$: \ $\res{r:\ x=1,\ y=-1+\lambda,\ z=\lambda}$

\ex{1} $\overrightarrow{AB}=(1,\,1,\,-2)$, $\overrightarrow{AC}=(-1,\,3,\,-1)$: \ $\left\{\begin{aligned}x&=1+\lambda-\mu\\y&=\lambda+3\mu\\z&=2-2\lambda-\mu\end{aligned}\right.$; \ $\vec n=\overrightarrow{AB}\times\overrightarrow{AC}=(5,\,3,\,4)$: \ $\res{5x+3y+4z-13=0}$

\ex{2} $\vec n=(2,\,1,\,-1)\times(1,\,0,\,3)=(3,\,-7,\,-1)$: \ $3(x-2)-7(y+1)-(z-1)=0$, \ $\res{3x-7y-z-12=0}$

\ex{2} $\vec d_s=(-2,\,-4,\,2)=-2\vec d_r$ y $Q(-1,\,2,\,1)\notin r$ ($\tfrac{-2}{1}\neq\tfrac32$): paralelas. Plano por $P(1,\,-1,\,0)$ con vectores $\vec d_r$ y $\overrightarrow{PQ}=(-2,\,3,\,1)$: \ $\vec n=(5,\,1,\,7)$, \ $\res{5x+y+7z-4=0}$

\ex{2} \mbox{$\dfrac x2+\dfrac{y}{-3}+\dfrac z6=1$}, \ $\res{3x-2y+z-6=0}$. \ El plano $4x-2y+3z=12$ corta los ejes en $(3,\,0,\,0)$, $(0,\,-6,\,0)$ y $(0,\,0,\,4)$: \ \mbox{$V=\dfrac{3\cdot6\cdot4}{6}=\res{12\ u^3}$}

\ex{2} $\vec d_r=(1,\,1,\,0)\times(0,\,1,\,-1)=(-1,\,1,\,1)$, que es normal al plano: \ $\res{x-y-z=0}$. \ La recta tiene vector director $(2,\,-1,\,3)$: \ $\res{(x,\,y,\,z)=(1,\,2,\,3)+\lambda(2,\,-1,\,3)}$

\ex{1} $\vec n=(1,\,2,\,0)\times(0,\,1,\,3)=(6,\,-3,\,1)$: \ $\res{6x-3y+z-5=0}$. \ $\tfrac{1+1}{2}=1$: \ $a=2+3=5$, $b=1-1=0$: \ $\res{a=5,\ b=0}$

\ex{1} \pau{Andalucía, ordinaria 2025}Plano por $P$ con vectores $\vec d_r=(1,\,2,\,2)$ y $\overrightarrow{QP}=(0,\,-1,\,-2)$, donde $Q(1,\,2,\,3)\in r$: \ $\vec n=(-2,\,2,\,-1)$: \ $\res{2x-2y+z-1=0}$

\ex{2} \pau{Aragón, ordinaria 2025}\begin{tasks}(1)
  \task $\res{x=4+\lambda,\ y=-3-2\lambda,\ z=\lambda}$
  \task $\vec d_r=(1,\,1,\,1)\times(1,\,-2,\,0)=(2,\,1,\,-3)$, $R(0,\,0,\,1)\in r$, $\overrightarrow{RQ}=(1,\,2,\,2)$: \ $\vec n=(8,\,-7,\,3)$: \ $\res{8x-7y+3z-3=0}$
\end{tasks}

\ex{2} \pau{Cantabria, extraordinaria 2024}\begin{tasks}(1)
  \task $\res{(x,\,y,\,z)=(0,\,3,\,2)+\alpha(4,\,-2,\,1)+\beta(2,\,0,\,2)}$
  \task $\big[\overrightarrow{AB},\overrightarrow{AC},\overrightarrow{AD}\big]=\begin{vmatrix}4 & -2 & 1\\2 & 0 & 2\\0 & -2 & 0\end{vmatrix}=12$: \ $V=\res{2\ u^3}$
\end{tasks}

\ex{2} \pau{Castilla-La Mancha, extraordinaria 2026}\begin{tasks}(1)
  \task $\overrightarrow{AB}\times\overrightarrow{AC}=(1,\,-1,\,1)\times(-1,\,1,\,1)=(-2,\,-2,\,0)$: \ $\res{x+y-2=0}$
  \task $\big[\overrightarrow{AB},\overrightarrow{AC},\overrightarrow{AD}\big]=(-2,\,-2,\,0)\cdot(a-1,\,0,\,2)=2-2a$; \ $\tfrac{|2-2a|}{6}=3$: \ $\res{a=-8\ \text{o}\ a=10}$
\end{tasks}

\ex{2} \pau{Castilla y León, extraordinaria 2026}\begin{tasks}(1)
  \task $E$ es el punto medio de las diagonales: $C=2E-A=\res{(-1,\,-1,\,1)}$, \ $D=2E-B=\res{(0,\,-2,\,2)}$
  \task $\overrightarrow{AB}\times\overrightarrow{AE}=(-1,\,1,\,-1)\times(-1,\,-1,\,0)=(-1,\,1,\,2)$: \ $\res{x-y-2z+2=0}$
  \task $\big[\overrightarrow{OA},\overrightarrow{OB},\overrightarrow{OE}\big]=\begin{vmatrix}1 & 1 & 1\\0 & 2 & 0\\0 & 0 & 1\end{vmatrix}=2$: \ $V=\res{\tfrac13\ u^3}$
\end{tasks}

\ex{2} \pau{Cataluña, ordinaria 2026}\begin{tasks}(1)
  \task $\overrightarrow{PQ}\times\overrightarrow{PR}=(2,\,-2,\,1)\times(0,\,1,\,2)=(-5,\,-4,\,2)$: \ $\res{5x+4y-2z-7=0}$
  \task $\tfrac12|(-5,\,-4,\,2)|=\tfrac12\sqrt{45}=\tfrac{3\sqrt5}{2}$ \ \checkmark
  \task $\tfrac16\big|(-5,\,-4,\,2)\cdot(x-1,\,y,\,z+1)\big|=1\iff|5x+4y-2z-7|=6$: $S$ debe estar en uno de los planos \ $\res{5x+4y-2z-1=0}$ \ o \ $\res{5x+4y-2z-13=0}$
\end{tasks}

\ex{2} \pau{Extremadura, extraordinaria 2026}\begin{tasks}(1)
  \task $\big[\overrightarrow{AB},\overrightarrow{AC},\overrightarrow{AD}\big]=\begin{vmatrix}-1 & 0 & 1\\-3 & -3 & 1\\0 & -3 & -2\end{vmatrix}=0$: son coplanarios. \ $\vec n=(3,\,-2,\,3)$: \ $\res{3x-2y+3z-15=0}$
  \task $r:(4-3t,\ 3-3t,\ 3+t)$; $z=2$ da $t=-1$ y $x=7$ \checkmark: \ $\res{k=6}$
\end{tasks}

\ex{2} \pau{Extremadura, ordinaria 2026}\begin{tasks}(1)
  \task $\vec u\cdot\vec v=0+12-12=0$: rectángulo. \ $\vec u\times\vec v=(25,\,-48,\,-36)$, de módulo $\sqrt{4225}$: \ área $=\res{65\ u^2}$
  \task $\res{25x-48y-36z-25=0}$
  \task $\pm\tfrac{13}{65}(25,\,-48,\,-36)=\res{\pm\left(5,\,-\tfrac{48}{5},\,-\tfrac{36}{5}\right)}$
\end{tasks}

\ex{2} \pau{Madrid, ordinaria 2025}Vectores directores $\vec d_r=(2,\,0,\,1)$ y $\vec n_\pi=(1,\,2,\,-3)$, punto $(1,\,0,\,2)$: \ $\vec n=(2,\,0,\,1)\times(1,\,2,\,-3)=(-2,\,7,\,4)$: \ $\res{2x-7y-4z+6=0}$


% =====================================================
\seccio{6}{Posiciones relativas de rectas y planos}

\ex{1} $\vec d_s=2\vec d_r$ y $Q(3,\,-3,\,4)\notin r$ ($\tfrac22=1$ pero $\tfrac43\neq1$): \ \rest{paralelas}. Plano por $P(1,\,-2,\,0)$ con $\vec d_r$ y $\overrightarrow{PQ}=(2,\,-1,\,4)$: \ $\res{x+2y+3=0}$

\ex{1} Los vectores directores no son proporcionales y $\big[\vec d_r,\vec d_s,\overrightarrow{PQ}\big]=\begin{vmatrix}2 & 1 & -1\\1 & -2 & 1\\-1 & 3 & -1\end{vmatrix}=-3\neq0$: \ \rest{se cruzan}.

\ex{2} $\overrightarrow{PQ}=(-1,\,-1,\,3-a)$: \ $\big[\vec d_r,\vec d_s,\overrightarrow{PQ}\big]=\begin{vmatrix}1 & 2 & 1\\1 & 0 & 1\\-1 & -1 & 3-a\end{vmatrix}=2a-8=0\Rightarrow\res{a=4}$. \ $1+\lambda=\mu$, $2+2\lambda=1$, $4+\lambda=3+\mu$ $\Rightarrow$ $\lambda=-\tfrac12$, $\mu=\tfrac12$: \ $\res{\left(\tfrac12,\,1,\,\tfrac72\right)}$

\ex{2} $\vec n\cdot\vec d=3-1-2=0$: la recta es paralela al plano o está contenida en él. $P(2,\,-1,\,0)$: $6+1+k=0\iff k=-7$. \ \rest{Si $k=-7$, $r\subset\pi$; si $k\neq-7$, paralelos.}

\ex{2} $r$: $(1,\,\lambda,\,\lambda)$. \ $\vec n\cdot\vec d=(1,\,m,\,1)\cdot(0,\,1,\,1)=m+1$. \ \rest{Si $m\neq-1$, secantes}; si $m=-1$, $P(1,\,0,\,0)$ da $1\neq2$: \rest{paralelos}. \ Para $m=1$: $1+2\lambda=2$, $\lambda=\tfrac12$: \ $\res{\left(1,\,\tfrac12,\,\tfrac12\right)}$

\ex{2} Si $a\neq-2$: \rest{secantes}. Si $a=-2$, $b=2$: \rest{coincidentes}. Si $a=-2$, $b\neq2$: \rest{paralelos}. \ Para $a=1$, $b=2$: $\vec d=(2,\,-1,\,3)\times(4,\,1,\,6)=(-9,\,0,\,6)\parallel(3,\,0,\,-2)$ y, con $z=0$, $x=\tfrac12$, $y=0$: \ $\res{(x,\,y,\,z)=\left(\tfrac12,\,0,\,0\right)+\lambda(3,\,0,\,-2)}$

\ex{3} Plano por $P$ paralelo a $\pi$: $x+y+z=3$. Corta a $s$: $2+1+\mu+\mu=3$, $\mu=0$: $X(2,\,1,\,0)$. \ $\overrightarrow{PX}=(1,\,0,\,-1)$: \ $\res{(x,\,y,\,z)=(1,\,1,\,1)+\lambda(1,\,0,\,-1)}$

\ex{2} \pau{Asturias, extraordinaria 2024}\begin{tasks}(1)
  \task $\big[\overrightarrow{AB},\overrightarrow{AC},\overrightarrow{AD}\big]=\begin{vmatrix}0 & -1 & 1\\-2 & 0 & 2\\-2 & -1 & 0\end{vmatrix}=6\neq0$: \ \rest{no hay ninguno}.
  \task $\overrightarrow{AB}\times\overrightarrow{AC}=(-2,\,-2,\,-2)$: $\pi:x+y+z-3=0$; \ \mbox{$\res{r:\dfrac{x+1}{1}=\dfrac y1=\dfrac{z-1}{1}}$}
  \task $(-1+t)+t+(1+t)-3=0$, $t=1$: \ $\res{P(0,\,1,\,2)}$
\end{tasks}

\ex{2} \pau{Canarias, extraordinaria 2025}\begin{tasks}(1)
  \task $\vec d_r=(1,\,1,\,-2)\times(2,\,1,\,-6)=(-4,\,2,\,-1)$, $R(-1,\,4,\,0)\in r$; $\vec d_s=(1,\,3,\,2)$, $Q(1,\,0,\,-1)$: \ $\big[\vec d_r,\vec d_s,\overrightarrow{RQ}\big]=\begin{vmatrix}-4 & 2 & -1\\1 & 3 & 2\\2 & -4 & -1\end{vmatrix}=0$ \ \checkmark
  \task $\vec d_r\times\vec d_s=(7,\,7,\,-14)\parallel(1,\,1,\,-2)$, por $Q$: \ $\res{x+y-2z-3=0}$
\end{tasks}

\ex{2} \pau{Canarias, ordinaria 2026}\begin{tasks}(1)
  \task Vector director $\vec n_\pi=(2,\,3,\,0)$: \ $\res{r:\ x=1+2\lambda,\ y=-5+3\lambda,\ z=1}$
  \task $\vec d_s=(-1,\,3,\,-3)\times(4,\,9,\,-2)=(21,\,-14,\,-21)\parallel(3,\,-2,\,-3)$; \ $\vec n\cdot\vec d_s=6-6=0$ y el punto $\left(1,\,-\tfrac23,\,0\right)\in s$ cumple $2-2=0$: \ \rest{$s$ está contenida en $\pi$}.
\end{tasks}

\ex{3} \pau{Canarias, extraordinaria 2024}\begin{tasks}(1)
  \task $\vec d_1=(7,\,7,\,7)\parallel(1,\,1,\,1)$, $P_1(1,\,1,\,0)$; \ $\vec d_2=(-2,\,1,\,-2)$, $P_2(1,\,0,\,1)$: \ $\big[\vec d_1,\vec d_2,\overrightarrow{P_1P_2}\big]=\begin{vmatrix}1 & 1 & 1\\-2 & 1 & -2\\0 & -1 & 1\end{vmatrix}=3\neq0$: \ \rest{se cruzan}.
  \task $\vec d_1\times\vec d_2=(-3,\,0,\,3)\parallel(1,\,0,\,-1)$: \ $s:\left(\lambda,\,\tfrac12,\,-\lambda\right)$. Con $r_1:(1+\mu,\,1+\mu,\,\mu)$: $\mu=-\tfrac12$, $\lambda=\tfrac12$: \ $\res{\left(\tfrac12,\,\tfrac12,\,-\tfrac12\right)}$
\end{tasks}

\ex{2} \pau{Cantabria, ordinaria 2024}$\vec d_r=(1,\,1,\,1)\times(1,\,2,\,-1)=(-3,\,2,\,1)$, $R(-10,\,5,\,0)\in r$; \ $\vec n\cdot\vec d=-6+2-a=-4-a$. Si $a=-4$, $R$ da $-15\neq3$.\begin{tasks}(1)
  \task \rest{Ningún valor de $a$.}
  \task \rest{$a=-4$.}
  \task \rest{Cualquier $a\neq-4$}; por ejemplo, con $a=2$ se cortan en $\res{(-1,\,-1,\,-3)}$.
\end{tasks}

\ex{2} \pau{Castilla y León, ordinaria 2025}\begin{tasks}(1)
  \task $P(0,\,1,\,0)\in r$, $Q(1,\,m,\,3)$: $\overrightarrow{PQ}=(1,\,m-1,\,3)=\vec d_s$, así que $\big[\vec d_r,\vec d_s,\overrightarrow{PQ}\big]=0$, y $\vec d_r=(4,\,-2,\,2)$ no es paralelo a $\vec d_s$: \rest{secantes para todo $m$}. (De hecho, $Q-\vec d_s=P$: el punto $P$ es de $s$.)
  \task $\res{(0,\,1,\,0)}$ (para cualquier $m$).
\end{tasks}

\ex{2} \pau{Castilla y León, extraordinaria 2025}\begin{tasks}(1)
  \task $\vec n\cdot\vec d=-5-m+2=-3-m$; $P(1,\,1,\,0)$: $5+m=2\iff m=-3$. Si $m=-3$: \rest{$r\subset\pi$}; si $m\neq-3$: \rest{secantes}.
  \task $\vec n'=(-1,\,-1,\,2)\times(5,\,1,\,1)=(-3,\,11,\,4)$: \ $\res{3x-11y-4z+8=0}$
\end{tasks}

\ex{1} \pau{Galicia, extraordinaria 2025}\begin{tasks}(1)
  \task $\vec d_s=2\vec d_r$ y $Q(2,\,1,\,1)\notin r$: \ \rest{paralelas}.
  \task $\vec n=(2,\,-1,\,1)\times\overrightarrow{PQ}=(2,\,-1,\,1)\times(1,\,-1,\,0)=(1,\,1,\,-1)$: \ $\res{x+y-z-2=0}$
\end{tasks}

\ex{2} \pau{Galicia, extraordinaria 2024}\begin{tasks}(1)
  \task $\vec d=(-2,\,0,\,-1)$: $\vec n\cdot\vec d=-2+2=0$ y $A$: $2+2-4=0$: \ \rest{$r$ está contenida en $\pi$}.
  \task $\vec n'=(-2,\,0,\,-1)\times(1,\,2,\,-2)=(2,\,-5,\,-4)$: \ $\res{2x-5y-4z+9=0}$
\end{tasks}

\ex{1} \pau{Galicia, ordinaria 2024}$\vec n\cdot\vec d=12+4+4b=0\Rightarrow b=-4$; el punto $(2,\,c,\,3)$ debe ser de $\pi$: $8+2c-12=2$: \ $\res{b=-4,\ c=3}$

\ex{2} \pau{Murcia, extraordinaria 2025}\begin{tasks}(1)
  \task $\vec d=(3,\,-2,\,-1)$, $R(2,\,-1,\,0)\in r$: $\vec n\cdot\vec d=4-a=0$ y $4-1=b$: \ $\res{a=4,\ b=3}$
  \task \rest{Si $a\neq4$} (cualquier $b$). Con $a=0$, $b=7$: $2(2+3t)+(-1-2t)=7$, $t=1$: \ $\res{(5,\,-3,\,-1)}$
\end{tasks}

\ex{2} \pau{Navarra, extraordinaria 2024}\begin{tasks}(1)
  \task $r:(-3+2t,\ 7-3t,\ t)$; en $\pi$: $2t-4=0$, $C(1,\,1,\,-2)$. \ $\overrightarrow{AB}\times\overrightarrow{AC}=(-2,\,-1,\,0)\times(-2,\,-1,\,-1)=(1,\,-2,\,0)\neq\vec 0$: \ \rest{no alineados}.
  \task $\tfrac12|(1,\,-2,\,0)|=\res{\tfrac{\sqrt5}{2}\ u^2}$
\end{tasks}

\ex{3} \pau{Navarra, ordinaria 2025}Plano por $P$ que contiene a $s$ (haz): $(2x+y-3z-6)+\lambda(2x-3z-8)=0$; en $P$: $1-\lambda=0$: $4x+y-6z-14=0$. \ Plano por $P$ que contiene a $r$: $\vec n=(2,\,1,\,1)\times(3,\,0,\,1)=(1,\,1,\,-3)$: $x+y-3z-5=0$. \ $\vec d_t=(4,\,1,\,-6)\times(1,\,1,\,-3)=(3,\,6,\,3)$: \ \mbox{$\res{t:\dfrac{x-2}{1}=\dfrac y2=\dfrac{z+1}{1}}$} \ (corta a $r$ en $(3,\,2,\,0)$ y a $s$ en $(1,\,-2,\,-2)$).

\ex{2} \pau{País Vasco, extraordinaria 2024}\begin{tasks}(1)
  \task $\vec d=(2,\,-1,\,1)\times(1,\,-1,\,4)=(-3,\,-7,\,-1)$: $\vec n\cdot\vec d=-6+21-A=0$, $A=15$; y $(-1,\,-2,\,0)\in r$ da $4\neq10$: \ $\res{A=15}$
  \task $\res{(2,\,5,\,1)}$
\end{tasks}

\ex{1} \pau{Aragón, ordinaria 2026}\begin{tasks}(1)
  \task Los vectores normales $(1,\,-2,\,4)$ y $(2,\,-3,\,6)$ no son proporcionales: \ \rest{no}: siempre se cortan en una recta.
  \task $\vec d=(0,\,2,\,1)$ y, con $z=0$, $(-1,\,0,\,0)$: \ $\res{x=-1,\ y=2\lambda,\ z=\lambda}$
\end{tasks}


% =====================================================
\seccio{7}{Posición relativa de tres planos}

\ex{2} $|M|=0$ y $\mathrm{rg}\,M=2$ (de hecho, la 3.ª fila de $M$ es la 2.ª menos la 1.ª). $\mathrm{rg}\,M^*=2\iff a=2-1=1$. \ \rest{Si $a=1$: secantes en una recta; si $a\neq1$: prisma} (no hay dos planos paralelos). \ Recta: \ $\res{x=1,\ y=\lambda,\ z=\lambda}$

\ex{3} $|M|=(k-1)^2(k+2)$. \ \rest{Si $k\neq1,-2$: secantes en un punto.} \ \rest{Si $k=1$: coincidentes} (los tres son $x+y+z=1$). \ Si $k=-2$: $\mathrm{rg}\,M=2$ y, sumando las tres ecuaciones, $0=-6$: $\mathrm{rg}\,M^*=3$ y ningún par es paralelo: \rest{prisma}. \ Punto: por simetría $x=y=z$, y $(k+2)x=k$: \ $\res{x=y=z=\tfrac{k}{k+2}}$

\ex{2} Hace falta que el tercero sea combinación de los dos primeros: $2\pi_1+\pi_2$ da $3x+y+4z=5$: \ $\res{a=4,\ b=5}$. \ Si $a=4$ y $b\neq5$: $\mathrm{rg}\,M=2$, $\mathrm{rg}\,M^*=3$: \rest{prisma}. \ Si $a\neq4$: \rest{secantes en un punto}.

\ex{2} \begin{tasks}(1)
  \task Debe ser paralelo a la recta $\pi_1\cap\pi_2$ (vector director $(1,\,1,\,1)$) sin contenerla; por ejemplo, $\res{x-z=1}$.
  \task Cualquier plano no paralelo a $(1,\,1,\,1)$; por ejemplo, $\res{z=0}$.
\end{tasks}

\ex{3} \pau{La Rioja, extraordinaria 2026}$|M|=\begin{vmatrix}1 & 1 & 1\\a-1 & 1 & 1\\1 & a-1 & -1\end{vmatrix}=a(a-2)$.\begin{tasks}(1)
  \task \rest{Si $a\neq0$ y $a\neq2$.} Restando $\pi_1$ de $\pi_2$: $(a-2)x=a-1$; \ $\res{\left(\tfrac{a-1}{a-2},\ \tfrac{1}{2-a},\ 0\right)}$
  \task Si $a=0$: $\pi_2:-x+y+z=0$ y $\pi_3:x-y-z=0$ son coincidentes y cortan a $\pi_1$: \ \rest{$a=0$}, \ $\res{x=\tfrac12,\ y=\tfrac12-\lambda,\ z=\lambda}$. \ (Si $a=2$, $\pi_1\parallel\pi_2$: no hay ningún punto común.)
\end{tasks}


% =====================================================
\seccio{8}{Haces de planos}

\ex{1} $(x-y+z-1)+\lambda(2x+y-3)=0$. \ En $O$: $-1-3\lambda=0$, $\lambda=-\tfrac13$: \ $\res{x-4y+3z=0}$

\ex{2} $\vec n=(1+2\lambda,\ -1+\lambda,\ 1)$. Paralelo a $OZ$ quiere decir $\vec n\cdot(0,\,0,\,1)=0$, es decir, $1=0$: ¡ningún $\lambda$! El plano que perdemos en el haz, $2x+y-3=0$, sí es paralelo a $OZ$: \ $\res{2x+y-3=0}$. \ Perpendicular a $\pi$: $\vec n\cdot(1,\,2,\,-1)=4\lambda-2=0$, $\lambda=\tfrac12$: \ $\res{4x-y+2z-5=0}$

\ex{2} $2x-y+2z+k=0$, \ $2+1+4+k=0$: \ $\res{2x-y+2z-7=0}$. \ El plano $2x-y+2z=k$ corta los ejes en $\left(\tfrac k2,0,0\right)$, $(0,\,-k,\,0)$, $\left(0,0,\tfrac k2\right)$: \mbox{$V=\dfrac{|k|^3}{24}=9\Rightarrow|k|=6$}: \ $\res{2x-y+2z=\pm6}$

\ex{2} $r:\left\{\begin{aligned}x-2y-1&=0\\3y-z-1&=0\end{aligned}\right.$; \ haz $(x-2y-1)+\lambda(3y-z-1)=0$; \ en $P$: $-2+\lambda=0$, $\lambda=2$: \ $\res{x+4y-2z-3=0}$

\ex{2} \pau{Murcia, extraordinaria 2026}Haz: $(2x+y-z-4)+\lambda(x+y-2z-3)=0$; en $P$: $-4+\lambda=0$, $\lambda=4$: \ $\res{6x+5y-9z-16=0}$

\ex{2} \pau{País Vasco, ordinaria 2025}$r_2\cap\pi$: $3t-3(1-2t)-2(2+t)+7=7t=0$, punto $(0,\,1,\,2)$. \ Haz: $(x+y-2z)+\lambda(2x-3y+z-1)=0$; en $(0,\,1,\,2)$: $-3-2\lambda=0$, $\lambda=-\tfrac32$: \ $\res{4x-11y+7z-3=0}$


% =====================================================
\seccio{9}{Problemas y cuestiones}

\ex{2} \begin{tasks}(1)
  \task \rest{Cierta}: si ambos son no nulos, $\cos\alpha=0$ y $\sin\alpha=0$ a la vez, imposible.
  \task \rest{Cierta}: no pueden ser coplanarios (si $\vec w$ fuera combinación de $\vec u$ y $\vec v$, multiplicando escalarmente por $\vec w$ tendríamos $|\vec w|^2=0$).
  \task \rest{Falsa}: por ejemplo $(1,\,0,\,0)$, $(0,\,1,\,0)$, $(1,\,1,\,0)$.
  \task \rest{Falsa}: se pueden cruzar.
  \task \rest{Falsa}: basta con que $\vec u\times(\vec v-\vec w)=\vec 0$, es decir, $\vec v-\vec w\parallel\vec u$; por ejemplo $\vec u=\vec i$, $\vec v=\vec j$, $\vec w=\vec i+\vec j$.
\end{tasks}

\ex{2} Segmentaria: $\tfrac x2+\tfrac y2+\tfrac z2=1$: \ $\res{x+y+z=2}$. \ $\overrightarrow{AB}\times\overrightarrow{AC}=(4,\,4,\,4)$: área $=\tfrac12\cdot4\sqrt3=\res{2\sqrt3\ u^2}$. \ Diagonal $(t,\,t,\,t)$: $3t=2$: \ $\res{\left(\tfrac23,\,\tfrac23,\,\tfrac23\right)}$

\ex{2} Trayectoria: $(1+2\lambda,\ 2+\lambda,\ 3\lambda)$, $\lambda\ge0$. Tejado: $3+6\lambda=12$, $\lambda=\tfrac32$, punto $\left(4,\,\tfrac72,\,\tfrac92\right)$. Poste ($x=5$, $y=4$): $\lambda=2$ cumple ambas, punto $(5,\,4,\,6)$. \ \rest{Choca primero con el tejado} ($\tfrac32<2$).

\ex{2} \pau{Aragón, extraordinaria 2026}Las cuatro patas se apoyan si los cuatro extremos son coplanarios: \ $\big[\overrightarrow{AB},\overrightarrow{AC},\overrightarrow{AD}\big]=\begin{vmatrix}22 & 39 & 0\\1 & 1 & 45\\22 & 40 & 45\end{vmatrix}=-1755\neq0$: \ \rest{no son coplanarios: la silla cojeará}.

\ex{2} \pau{Islas Baleares, extraordinaria 2026}\begin{tasks}(1)
  \task $\big[\overrightarrow{OA},\overrightarrow{OB},\overrightarrow{OC}\big]=\begin{vmatrix}1 & 2 & 3\\3 & 2 & 1\\5 & 4 & 3\end{vmatrix}=0$: \ \rest{sí}.
  \task El plano $ODE$ tiene vector normal $(2,\,1,\,1)\times(4,\,3,\,1)=(-2,\,2,\,2)$: $x-y-z=0$. Cualquier punto de este plano (no alineado con $O$, $D$, $E$); por ejemplo, \ $\res{F(1,\,1,\,0)}$
\end{tasks}

\ex{2} \pau{Islas Baleares, ordinaria 2026}$\res{M_1(2,\,3,\,6),\ M_2(4,\,4,\,5),\ M_3(5,\,1,\,4)}$. \ $\overrightarrow{M_1M_2}\times\overrightarrow{M_1M_3}=(2,\,1,\,-1)\times(3,\,-2,\,-2)=(-4,\,1,\,-7)$: \ $\res{4x-y+7z-47=0}$

\ex{2} \pau{La Rioja, ordinaria 2026}\begin{tasks}(1)
  \task Vectores $\left(1,\,-2,\,\tfrac34\right)$ y $\left(-1,\,0,\,-\tfrac12\right)$; su producto vectorial es proporcional a $(4,\,-1,\,-8)$: \ $\res{4x-y-8z+64=0}$
  \task $4s-16-64+64=0$: \ $\res{s=4}$
\end{tasks}


\end{document}
